% ---
% id: protocol-compiler/def:gate-outcome-table
% kind: def
% title: "The gate outcome table"
% summary: >-
%   The gate outcome table gives each gate failure a refusal naming what failed
%   or an emission that is not key-length-eligible; its two certificate-absent
%   rows give the compiler no outcome. Where a declaration meets a refusal row
%   and an emission row, the refusal wins.
% prominence: standard
% status: implemented
% uses: [def:eligible]
% mentions: [thm:gates, req:gate-outcomes]
% assumes: []
% refs: []
% provenance: newly-derived
% implements: certiqs_key.protocol_compiler.gates:resolve_gates
% lean:
% ---

Rows 2 and 4 describe a case that does not reach the compiler: no
\tref{protocol specification} meets their premise (Theorem~\ref{thm:gates}\,(ii)), and the
transmitter refuses such a request before any artefact exists
(Requirement~\ref{req:gate-outcomes}). Their cell says so; the table gives
the compiler no outcome there.

\begin{definition}\label{def:gate-outcome-table}
The
\term[The table that assigns to each gate failure its outcome, a refusal or an emission that is not key-length-eligible, or none for the two certificate-absent rows, the refusal taking precedence where a declaration meets both kinds of row.]{gate outcome table}
is the following table of the failures of the gates of
Definition~\ref{def:eligible} and their outcomes; where the failures of one declaration meet the
premises of a refusal row and of an emission row both, the outcome is the
refusal.

\medskip
\noindent\emph{Gate failure to emitted outcome.} No row substitutes a default
value, drops an entry, changes the source mode or synthesizes an attestation.
\phantomsection\label{tab:gateoutcome}

\noindent{\small
\begin{tabular}{@{}>{\raggedright\arraybackslash}p{0.58\textwidth}
                  >{\raggedright\arraybackslash}p{0.36\textwidth}@{}}
\hline
\emph{What failed} & \emph{Outcome} \\
\hline
(G1): any of V1--V8 & refusal, naming the check \\
(G2): trusted mode at $d_A\ge2$, certificate \emph{absent}, and
$m=0$ &
none (no outcome of the compiler) \\
(G2): trusted mode \emph{at any $d_A$}, certificate \emph{present}
but one
of the five conditions of Definition~\ref{def:eligible}\,(G2) fails ---
$\Psi$ Hermitian, $\Psi^2=\Psi$, $\tr\Psi=1$ rounding nowhere, (SR1), (SR2) &
refusal, naming the condition \\
(G2): trusted mode at $d_A\ge2$, certificate absent, and $m\ge1$ &
none (no outcome of the compiler) \\
(G3)(a): $\CList$ not in the licensed class --- some entry declared
\tref{general} in \tref{trusted-source mode} &
emission, not key-length-eligible \\
(G3)(b): the attack-set licence record absent &
emission, not key-length-eligible \\
\hline
\end{tabular}}

\noindent The cell of rows 2 and 4 assigns the compiler no outcome.
\end{definition}
